\documentclass[11pt]{article}

% --- Standard, widely-available packages only (Overleaf / any TeX Live install) ---
\usepackage[utf8]{inputenc}
\usepackage[T1]{fontenc}
\usepackage{geometry}
\geometry{margin=1in}
\usepackage{graphicx}
\usepackage{booktabs}
\usepackage{array}
\usepackage{longtable}
\usepackage{enumitem}
\usepackage[colorlinks=true,linkcolor=blue,citecolor=blue,urlcolor=blue]{hyperref}
\usepackage{xcolor}
\usepackage{amsmath}

\newcommand{\code}[1]{\texttt{#1}}

\title{\textbf{Does Garbage Collection Disqualify Java from Agentic AI's Hot Path?}\\
\large An Empirical Study of Tail Latency Under Token-Streaming Workloads}

\author{Raghavendra Venkateshappa \\
\small Independent Researcher \\
\small Senior Staff Performance Engineer \\
\small \texttt{raghavjax@gmail.com}}
\date{}

\begin{document}

\maketitle

\begin{abstract}
A claim has taken hold in parts of the AI-infrastructure community: garbage-collected runtimes
are structurally unfit for the low, deterministic tail latency that real-time LLM token streaming
demands, and that performance-critical ``hot path'' workloads are consequently migrating to
non-garbage-collected systems languages. This claim rests largely on anecdote and on the
behavior of earlier JVM garbage collector generations, not on direct measurement against modern
low-pause collectors. We present an empirical study measuring P50, P99, and P99.9 latency of a
Java-based agent-orchestration service streaming tokens over Server-Sent Events, under three
garbage collectors --- G1, Generational ZGC, and Shenandoah --- across three heap sizes (512MB,
1GB, 2GB) and a concurrency sweep reaching 10{,}000 simultaneous streaming sessions, with GC pause
events correlated directly against observed latency spikes via JDK Flight Recorder (JFR).

We find the disqualification claim, taken as a blanket statement, does not hold: across all
three heap sizes, when concurrency stays within roughly half of what a given heap can sustain,
GC accounts for a small and stable share of latency spikes (typically under 5\%) and the three
collectors perform comparably. The claim is not simply false, however. At the edge of what a
given heap size can sustain, all three collectors show GC becoming a first-order factor in tail
latency and reliability --- and, more importantly, \textbf{they fail in three genuinely different
ways that a single ``use a modern collector'' recommendation obscures.} At a 2GB heap under
7{,}500 concurrent streams, G1 degraded gracefully with zero crashes across six independent runs;
Generational ZGC failed outright with complete out-of-memory termination, an outcome reproduced
identically across two independent test conditions; Shenandoah never crashed but exhibited the
highest and most variable tail latency of the three (P99.9 up to 1.7 seconds, substantial
run-to-run variance even under controlled conditions). We propose a plausible, though unproven,
mechanism for ZGC's earlier failure --- its colored-pointers/multi-mapping design carries
documented memory overhead beyond G1's or Shenandoah's. We additionally report a methodological
finding of independent interest: ambient CPU contention from unrelated processes on a shared,
non-dedicated benchmark host measurably distorts GC-attribution results at high concurrency, and
we demonstrate a targeted re-run protocol --- diagnosed via JFR's \code{jdk.CPULoad} event --- to
separate genuine collector behavior from measurement noise. We conclude that the right collector,
and the right memory headroom, depends on which failure mode a given system can tolerate least,
and that this must be tested against the system's actual workload shape rather than inferred
from general collector reputation.
\end{abstract}

\section{Introduction}

Java-based enterprise systems are increasingly asked to host agentic AI workloads: services that
call out to LLMs, stream token-by-token responses back to callers, and orchestrate multi-step
tool-calling pipelines --- all while retaining per-request conversation state in memory for the
duration of a session. This is a materially different allocation and latency profile from the
request/response workloads the JVM's garbage collectors were tuned against for two decades.

A recurring claim has followed this shift: that garbage collection makes the JVM structurally
unsuitable for the tight, predictable latency that live token streaming demands, and that teams
building this kind of infrastructure are better served by non-garbage-collected systems
languages. Versions of this claim circulate in engineering blog posts and community commentary,
sometimes citing GC pause behavior from collector generations (e.g., CMS, early G1) that are no
longer the default or even available in current JDK releases. The claim is rarely tested directly
against modern low-pause collectors --- Generational ZGC and Shenandoah --- under a workload
shaped like real token streaming.

This matters specifically for Java practitioners because the standard advice --- ``just use a
modern low-pause collector'' --- is repeated often enough that it functions as received wisdom.
If that advice is wrong, or right only within limits nobody has stated precisely, teams are
making collector and capacity-planning decisions on faith rather than evidence. This report tests
the claim directly: we built a dependency-free Java harness that simulates the specific
allocation and latency shape of an agent-orchestration service streaming LLM tokens, and we
measured what actually happens across three collectors, three heap sizes, and a wide concurrency
range.

\section{Related Work}

Existing coverage of this topic is thin and, where it exists, is not empirical in the sense this
report is. InfoQ's own \emph{Java Trends Report 2025} --- a survey by several members of InfoQ's
Java editorial team --- explicitly states that it does not go deep on garbage collection, leaving
the topic as an acknowledged gap in the publication's own coverage \cite{infoq-trends-2025}.
InfoQ's JVM-focused article section's most substantial garbage-collection deep dive,
\emph{Understanding Classic Java Garbage Collection} \cite{evans-2020}, predates Generational
ZGC's production availability entirely. A closely related talk --- \emph{Java Concurrency from
the Trenches: Lessons Learned in the Wild} (Netflix, InfoQ Dev Summit Boston 2025)
\cite{marques-2025} --- documents a real production migration through parallel streams,
executor-based concurrency, and virtual threads for a high-throughput gRPC workload, including
an encounter with virtual-thread carrier pinning. It is a genuine practitioner account of
concurrency-model tuning, but it is a single company's architecture narrative, not a controlled
comparison of garbage collectors, and does not address the disqualification claim directly.

Beyond InfoQ, informal industry commentary --- blog posts and vendor-adjacent analyses --- has at
times framed the token-streaming ``hot path'' as migrating toward non-garbage-collected systems
languages on the premise that GC pauses are incompatible with streaming latency requirements. We
were unable to find a source that tested this premise directly against modern low-pause
collectors under a workload shaped like real token streaming, with reproducible methodology and
public data. That is the gap this report addresses.

\section{Methodology}

\subsection{Harness design}

The benchmark harness (\code{gc-agentic-benchmark/}) is a from-scratch Java project with zero
external dependencies --- it uses only \code{java.net.http}, \code{com.sun.net.httpserver}, and
\code{jdk.jfr.consumer}, all part of the standard JDK. This was a deliberate choice: it keeps
observed GC and latency behavior attributable to the JVM and collector under test, not to the
overhead of a specific web framework or reactive library layered on top.

\begin{itemize}[leftmargin=1.5em]
  \item \textbf{\code{StreamingServer}} --- an SSE (Server-Sent Events) endpoint that streams a
    configurable number of tokens per request, each with a delay drawn from
    \code{SyntheticTokenSource} (\S\ref{sec:workload}). Every request maintains a
    \code{ConversationContext} that accumulates the streamed tokens into a growing buffer and
    periodically snapshots it into a retained history list --- modeling the specific allocation
    and retention pattern of an agent-orchestration layer holding growing conversation state, not
    a synthetic no-op loop. This is what makes the workload actually exercise GC and heap
    capacity, rather than merely exercising I/O.
  \item \textbf{\code{LoadGenerator}} --- opens $N$ concurrent SSE streams, one Java virtual
    thread per stream, each issuing a configurable number of sequential requests and recording
    per-token arrival latency to a log file via a dedicated, bounded-queue writer thread
    (\S\ref{sec:bugs} discusses a bug found in this component).
  \item \textbf{\code{LatencyGcCorrelator}} --- reads each run's JFR recording, extracts
    \code{jdk.GarbageCollection} events (portable across all three collectors), computes latency
    percentiles from the recorded per-token gaps, and reports what fraction of latency spikes
    temporally overlap a GC pause event (\S\ref{sec:attribution} details the precise
    spike-attribution methodology and its limits).
\end{itemize}

Both server and client run as virtual threads (\code{Executors.newVirtualThreadPerTaskExecutor()}),
which is what makes opening thousands of concurrent, long-lived streaming sessions from a single
process practical --- a traditional one-OS-thread-per-connection model would not scale to the
concurrency levels tested here without the memory and scheduling cost of thousands of live OS
threads.

\subsection{Workload shape}
\label{sec:workload}

Each simulated session requests 1{,}000 tokens over SSE. Token timing is synthetic, generated by
\code{SyntheticTokenSource}: a time-to-first-token delay drawn from a Gaussian distribution
(mean 200ms, stddev 60ms, floored at 1ms) followed by inter-token gaps drawn from a log-normal
distribution (median $\sim$20ms). This is a plausible shape for LLM streaming behavior, not a
calibration against any specific model provider's real traffic --- see \S\ref{sec:limitations}
(Limitations). Each concurrent virtual thread issues 3 sequential such sessions before completing.

\subsection{JVM, collectors, and a hardware discovery worth reporting}
\label{sec:jvm}

All experiments ran on JDK 21 (LTS), chosen deliberately over a newer non-LTS build also
available on the test machine (JDK 26), because JDK 21 or a comparable current LTS is what
InfoQ's actual readership is far more likely to run in production. All three collectors under
test --- G1, Generational ZGC (\code{-XX:+UseZGC -XX:+ZGenerational}), and Shenandoah
(\code{-XX:+UseShenandoahGC}) --- were run with default tuning beyond the collector-selection
flag itself; no attempt was made to hand-tune any collector's heuristics for this workload.

A methodological discovery is worth reporting here because it materially affects reproducibility
guidance for anyone attempting to replicate or extend this work on Apple Silicon hardware: the
JDK initially used for this study, installed via the Intel/Rosetta-prefixed Homebrew installation
common on machines that have been upgraded from Intel to Apple Silicon Macs over time, was an
\textbf{x86\_64 binary running under Rosetta 2 binary translation}, not a native arm64 build.
This was not discovered through a crash or an obvious error --- the harness ran and produced
plausible output --- but through a performance anomaly during analysis-tool development: a
simple string \code{indexOf} operation, benchmarked in isolation, took 45 seconds for 6 million
calls under the x86\_64/Rosetta JDK versus 0.45 seconds under a native arm64 build of the same
JDK version --- a 100x difference with no algorithmic explanation. Because the benchmark's server
and client processes were running under the same emulated JVM, this was not merely an
analysis-tool inconvenience; it was a direct confound on the latency measurements themselves. All
data reported in this study was collected (or, where affected, re-collected) under a native
arm64 JDK 21.0.11 build. Anyone reproducing this work on Apple Silicon should explicitly verify
JDK binary architecture (\code{file \$(which java)} should report \code{arm64}, not
\code{x86\_64}) before trusting timing-sensitive results.

\subsection{Spike-attribution methodology}
\label{sec:attribution}

For each run, we compute the median inter-token latency and define a ``spike'' as any inter-token
gap exceeding $3\times$ that run's own median --- a relative rather than absolute threshold,
chosen because an absolute millisecond value does not transfer across different token-rate
assumptions or timing conditions (a design choice validated during a calibration phase where an
absolute threshold produced inconsistent results across compressed and real-time runs). For each
spike, we check whether its time window overlaps any \code{jdk.GarbageCollection} event recorded
in the same run's JFR data, and report the fraction of spikes with such overlap as
``GC-attributable.''

This methodology is a \textbf{correlational heuristic, not a proof of causation.} Temporal
overlap between a latency spike and a GC event is consistent with the GC event causing the
spike, but does not rule out both being caused by a third factor (e.g., general CPU contention).
We encountered this limitation directly during interpretation of high-concurrency results
(\S\ref{sec:threemodes}) and disclose it explicitly rather than treat GC-attribution percentages
as a settled causal measure.

\subsection{Bugs found and fixed during this study}
\label{sec:bugs}

Five implementation bugs surfaced during the course of this work, each of which would have
silently produced incorrect results if uncaught. We report them because a study whose entire
premise is measuring GC behavior precisely has an obligation to disclose where its own
measurement apparatus initially failed, and because each is a generically useful lesson for
anyone building similar harnesses:

\begin{enumerate}[leftmargin=1.5em]
  \item \textbf{Zombie process contamination.} An early version of the server-launch script
    killed a wrapper shell process rather than the JVM it had spawned, orphaning that JVM. It
    continued listening on a shared port across multiple subsequent test runs, meaning several
    early calibration runs believed they were testing different collectors while actually hitting
    the same stale process throughout. Caught by cross-checking each run's actual JFR-recorded
    collector configuration against what the test intended --- a check now built permanently into
    every run via an automated \code{verify\_collector} step.
  \item \textbf{Quadratic-cost log parsing.} The JFR/latency correlation tool initially
    recompiled a regular expression on every field of every logged token --- inconsequential at
    low token counts, but over 60 million redundant compilations at the highest concurrency
    levels tested, which presented as an apparent hang rather than a clean failure. Fixed with
    hand-rolled parsing into primitive arrays.
  \item \textbf{Unbounded producer/consumer queue.} The load generator's latency logger used an
    unbounded queue between thousands of concurrent recording virtual threads and a single
    file-writer thread. Under extreme, artificially compressed timing (used only for fast
    calibration, never for the reported results), producers could outpace the writer badly enough
    to grow the queue to tens of millions of objects, causing severe GC pressure in the
    \emph{load generator itself} that again presented as an apparent hang. Fixed with a bounded
    queue and a blocking \code{put()}, which makes producers throttle naturally instead.
  \item \textbf{No timeout on either phase of an HTTP request.} When a server under test
    genuinely OOM-crashed mid-stream (an expected outcome at the edge of a heap's capacity, see
    \S\ref{sec:threemodes}), the client had no way to detect this: neither
    \code{HttpClient.send()} waiting for an initial response, nor the subsequent blocking read of
    a streaming response body, had any timeout configured. A single such crash silently stalled
    an entire multi-hour unattended test run for over an hour before being caught. Fixed with an
    explicit request-level timeout for the first phase and a scheduled watchdog that force-closes
    the response stream for the second.
  \item \textbf{Compressed-timing calibration underestimates real memory pressure.} A calibration
    phase used an artificial timing-compression factor (20x faster than real inter-token delays)
    to find interesting concurrency ranges quickly. This calibration predicted a 2GB heap would
    survive the full tested range up to 10{,}000 concurrent streams without incident. The real
    (uncompressed) study instead found a genuine capacity boundary around 7{,}500--10{,}000
    concurrent streams at that heap size (\S\ref{sec:threemodes}). We attribute the discrepancy
    to compressed timing altering virtual-thread scheduling contention enough to reduce the true
    peak concurrent memory overlap the calibration run actually achieved. We report this so that
    others using a similar compressed-calibration technique treat it as directional guidance
    only, not a precise predictor of capacity limits.
\end{enumerate}

\subsection{Controlling for ambient system load}
\label{sec:noise}

During interpretation of the results, we observed unusually large run-to-run variance in
GC-attribution percentages at the highest tested concurrency levels --- in one case, three
otherwise-identical repetitions of the same (heap, collector, concurrency) configuration produced
GC-attributable spike rates of 10.3\%, 71.8\%, and 72.8\%. Direct inspection of the corresponding
JFR recordings ruled out a GC-side explanation: both the high- and low-variance repetitions
showed the same GC event types and comparably small (sub-millisecond) individual pause durations,
and the lower-percentage repetition in fact contained \emph{more} GC events than the
higher-percentage one. We then examined JFR's \code{jdk.CPULoad} event, which reports total
machine CPU utilization alongside this JVM's own share of it, and found \code{machineTotal}
utilization at 85--100\% throughout these runs while this JVM's own share (\code{jvmUser}) was
frequently only single digits to $\sim$30\% --- a substantial share of the machine's CPU was
being consumed by unrelated processes (in this case, a video-conferencing application's
background helper process was independently found consuming over 600\% CPU --- i.e., more than
six full cores --- at the time of measurement).

This is disclosed as a limitation of running on a shared, non-dedicated development machine
rather than an isolated benchmark host (\S\ref{sec:limitations}), but we did not treat it as
simply unavoidable noise to be reported alongside the data uncritically. We quit the offending
applications and re-ran the affected cells (2GB heap, 5{,}000/7{,}500/10{,}000 concurrency, all
three collectors, 27 cells total) under materially reduced ambient load.
Section~\ref{sec:threemodes} reports both the original and the re-verified results and what
changed between them. Notably, the OOM/capacity-boundary findings proved largely independent of
this confound (OOM is a hard memory-exhaustion event, not directly a function of CPU scheduling),
while the fine-grained latency-variance findings were meaningfully affected --- this distinction
is discussed explicitly rather than glossed over.

\subsection{Experimental design}

Concurrency levels were swept asymmetrically per heap size, chosen from calibration to avoid
spending test time on configurations expected to fail outright from insufficient memory
(Table~\ref{tab:sweep}).

\begin{table}[h]
\centering
\caption{Concurrency levels tested per heap size.}
\label{tab:sweep}
\begin{tabular}{ll}
\toprule
Heap size & Concurrency levels tested \\
\midrule
512MB & 100, 500, 1000 \\
1GB & 100, 500, 1000, 2500 \\
2GB & 100, 500, 1000, 2500, 5000, 7500, 10000 \\
\bottomrule
\end{tabular}
\end{table}

A fourth heap tier (4GB) was scoped and partially calibrated but deliberately excluded from the
final study: it would have added an estimated 2--3.5 additional hours of runtime for a data point
whose only plausible contribution was confirming that the observed trend continues, not
answering the core research question, which the 512MB$\to$1GB$\to$2GB progression already does.

Every (heap, collector, concurrency) cell was run 3 times independently, and we report the mean
alongside the minimum and maximum across repetitions rather than the mean alone, specifically
because some cells (\S\ref{sec:threemodes}) show high enough variance that a mean alone would
misrepresent the data. In total, the reported dataset comprises 126 cells across the three heap
tiers, plus a 27-cell targeted re-verification pass at the highest concurrency levels
(\S\ref{sec:noise}).

\section{Results}

Full per-cell data (mean, minimum, and maximum across 3 repetitions) is available in
\code{data/results/aggregate.csv}. This section summarizes the shape of the findings.

\subsection{Convergence at moderate concurrency-to-memory ratios}

Across all three heap sizes, at concurrency levels within roughly half of what that heap size
can sustain, GC-attributable spike percentages are low and comparable across all three
collectors (Table~\ref{tab:convergence}).

\begin{table}[h]
\centering
\caption{GC-attributable spike rate (mean, range across 3 reps) at moderate concurrency.}
\label{tab:convergence}
\begin{tabular}{llll}
\toprule
Heap & Collector & Concurrency & GC-attributable \% \\
\midrule
512MB & G1 & 100 & 0.30\% (0.3--0.3\%) \\
512MB & ZGC & 100 & 2.77\% (2.7--2.9\%) \\
512MB & Shenandoah & 100 & 0.43\% (0.3--0.6\%) \\
1GB & G1 & 1000 & 4.13\% (4.1--4.2\%) \\
1GB & ZGC & 1000 & 3.20\% (2.6--3.8\%) \\
1GB & Shenandoah & 1000 & 1.10\% (1.0--1.2\%) \\
2GB & G1 & 2500 & 2.70\% (2.5--2.8\%) \\
2GB & ZGC & 2500 & 2.90\% (2.6--3.5\%) \\
2GB & Shenandoah & 2500 & 1.67\% (1.3--1.9\%) \\
\bottomrule
\end{tabular}
\end{table}

At this scale, the median (P50) latency itself remains pinned near the workload's own baseline
inter-token delay ($\sim$20ms), meaning GC is neither meaningfully inflating typical latency nor
producing a large share of tail spikes at these concurrency levels, for any of the three
collectors. \textbf{The blanket disqualification claim does not hold in this regime.}

\subsection{Divergence at the top of each heap's range}

Within each heap tier, GC-attributable percentages rise with concurrency, and rise faster as
concurrency approaches the heap's capacity limit. At 1GB and 2500 concurrency --- the top of that
tier's tested range --- Shenandoah (9.1\%) and ZGC (11.9\%) show markedly higher GC-attribution
than at lower concurrency, though G1 (4.5\%) rises more modestly. All three remain well short of
failure at this level, however; the sharpest divergence appears at the very top of the 2GB tier,
detailed next.

\begin{figure}[h]
\centering
\includegraphics[width=\textwidth]{figures/convergence-divergence.png}
\caption{GC-attributable spike rate vs.\ concurrency, by heap size and collector. Each panel is
one heap size; each line is one collector's mean GC-attributable spike rate across the
concurrency sweep. All three tiers show the same shape: flat and low through moderate
concurrency, then a sharp rise near that heap's capacity limit. The 2GB panel's dashed segment
and hollow marker for Generational ZGC show where it lost repetitions entirely (0/3 at 7{,}500;
2/3 at 10{,}000) --- a solid line would misrepresent that as smooth, uninterrupted behavior.}
\label{fig:convergence}
\end{figure}

\subsection{Three failure modes at the edge of 2GB's capacity}
\label{sec:threemodes}

At 5{,}000--10{,}000 concurrent streams against a 2GB heap, the three collectors diverge sharply
and in qualitatively different ways. We report both the original (ambient-load-affected) results
and the re-verified results from the targeted quiet-machine re-run (\S\ref{sec:noise}), because
the comparison itself is informative about which findings are robust (Table~\ref{tab:beforeafter}).

\begin{table}[h]
\centering
\caption{Before (uncontrolled ambient load) vs.\ after (quieted machine) at 2GB.}
\label{tab:beforeafter}
\small
\begin{tabular}{p{1.6cm}p{1.3cm}p{4.8cm}p{4.8cm}}
\toprule
Collector & Conc. & Before: outcome & After: outcome \\
\midrule
G1 & 7500 & 2/3 survived (1 OOM), 11.3--16.6\% & \textbf{3/3 survived, zero OOM}, 10.9--17.5\% \\
G1 & 10000 & 2/3 survived (1 OOM), 12.4--25.4\% & \textbf{3/3 survived, zero OOM}, 7.7--14.6\% \\
ZGC & 7500 & \textbf{0/3 survived (all OOM)} & \textbf{0/3 survived (all OOM)} --- identical \\
ZGC & 10000 & 1/3 survived, 44.3\% & 2/3 survived, 49.3--50.2\% \\
Shenandoah & 7500 & 10.3\%, 71.8\%, 72.8\% (extreme spread) & 26.2\%, 38.4\%, 39.3\% (still elevated, far tighter) \\
Shenandoah & 10000 & 7.4--15.1\% & 12.2--32.8\% \\
\bottomrule
\end{tabular}
\end{table}

Two conclusions follow directly from this comparison. First, \textbf{G1's apparent OOM failures
in the original run did not reproduce under controlled conditions} --- on the quieted machine, G1
survived the entire tested range up to 10{,}000 concurrent streams with zero crashes. This part
of the original finding was very likely an artifact of ambient CPU contention pushing an
already-marginal configuration over the edge, and we do not report G1 as failure-prone at this
heap size; the corrected finding is that G1 degrades gracefully throughout this range. Second,
\textbf{ZGC's complete failure at 7{,}500 concurrent streams reproduced identically} across both
runs --- zero of three repetitions survived under either ambient-load condition. This is treated
as a robust, reproducible finding: at a 2GB heap, ZGC exhausts available memory at a substantially
lower concurrency threshold than G1 does. ZGC's continued difficulty at 10{,}000 concurrency
(2 of 3 survived under quiet conditions, both at approximately 50\% GC-attributable spikes)
reinforces this. Third, Shenandoah's extreme variance shrank substantially under controlled
conditions but did not disappear --- it remains the collector with the widest spread of outcomes
at the top of this range, suggesting a real, if smaller than first observed, degree of intrinsic
run-to-run variability in how its concurrent collection cycles cope with sustained memory
pressure near a capacity boundary. Shenandoah is also the only one of the three collectors to
survive every single repetition at every tested concurrency level up to 10{,}000, including under
the original uncontrolled conditions --- it trades reliability against memory exhaustion for
volatility in tail latency (P99.9 reached 2.16 seconds in one 7{,}500-concurrency repetition).

\begin{figure}[h]
\centering
\includegraphics[width=0.75\textwidth]{figures/failure-modes-7500.png}
\caption{Three failure modes at the edge of capacity: same heap, same concurrency, three
different outcomes. G1 degrades and keeps running. Generational ZGC exhausts memory outright ---
zero of three repetitions completed. Shenandoah completes every repetition but with the highest
mean and widest spread of GC-attributable spikes of the three. Whiskers show min--max across the
3 repetitions from the re-verified (quiet-machine) run.}
\label{fig:failuremodes}
\end{figure}

\section{Discussion}

The evidence supports a more specific claim than either ``GC disqualifies Java'' or its blanket
rebuttal. Within a reasonable concurrency-to-memory ratio, all three modern collectors handle
this workload comparably well, and GC is a minor contributor to tail latency --- the
disqualification claim, stated without qualification, is not supported by this data. But the
claim contains a displaced kernel of truth: push any of these collectors close enough to a given
heap's actual capacity, and GC absolutely becomes a first-order reliability and latency concern.
The more actionable and less obvious finding is that \textbf{this breakdown does not look the
same across collectors}, which means the common advice to simply select a modern low-pause
collector is itself an oversimplification. A team whose priority is avoiding hard crashes at
almost any latency cost might reasonably prefer Shenandoah's behavior in this data; a team that
can tolerate occasional degraded throughput but needs predictable failure behavior might prefer
G1's more graceful degradation; a team running ZGC at a heap size this close to its true
working-set requirement should specifically re-examine whether it has provisioned enough
headroom, given that ZGC was the first and most completely to fail in this data.

We offer one plausible, mechanistically grounded explanation for why ZGC failed earliest at a
fixed heap size, without claiming it as proven from this data alone: ZGC's design relies on
colored pointers and multi-mapped memory, which carries documented metadata and address-space
overhead beyond what G1 or Shenandoah require for the same nominal heap size. If that overhead is
the operative mechanism, the practical implication is straightforward --- ZGC should be
provisioned with more headroom relative to a measured working set than G1 or Shenandoah would
need for the same workload, not the same headroom on the assumption that ``modern collector''
implies ``equivalent or lower memory requirement.''

A secondary contribution of this work is methodological: we demonstrate that ambient load on a
shared benchmark host is not merely a hypothetical threat to validity but one that measurably
altered both the specific GC-attribution percentages and at least one apparent failure boundary
in this exact dataset, and we show a concrete, cheap way to detect it (JFR's \code{jdk.CPULoad}
event) and to control for it (a targeted, quieted re-run of only the affected cells, verified
rather than assumed to be sufficient). This is offered as a reusable technique for anyone
benchmarking JVM performance on infrastructure they do not fully control.

\section{Limitations}
\label{sec:limitations}

\begin{itemize}[leftmargin=1.5em]
  \item \textbf{Token timing is synthetic}, not calibrated against a specific model provider's
    real traffic. It is a plausible shape (a longer initial delay followed by shorter,
    log-normally distributed inter-token gaps), not a measured fit to any production LLM's actual
    streaming behavior.
  \item \textbf{Single, non-dedicated hardware.} All results were collected on one Apple Silicon
    development laptop. Section~\ref{sec:noise} shows this measurably affected fine-grained
    results at the highest concurrency levels even after a controlled re-run; a fully isolated
    benchmark host (bare-metal or a dedicated cloud instance with no co-tenants) was not used and
    would be a natural next step for anyone extending this work.
  \item \textbf{A single request/allocation shape.} All results use a fixed shape --- 1{,}000
    tokens per session, 3 sequential sessions per concurrent stream, and a fixed
    conversation-context retention pattern. The specific concurrency levels at which each
    collector's capacity boundary appears are almost certainly shape-dependent; the qualitative
    finding (collectors converge at moderate load, diverge differently near capacity) is expected
    to generalize better than the precise numeric thresholds are.
  \item \textbf{Default GC tuning only.} No collector-specific tuning flags (e.g., ZGC's
    soft-max-heap setting, Shenandoah's pacing heuristics, G1's region sizing) were explored. It
    is plausible that targeted tuning would shift or narrow the gaps reported here, particularly
    for ZGC's early failure.
  \item \textbf{Compressed-timing calibration was shown, in this same study, to underestimate
    real memory pressure at high concurrency} (\S\ref{sec:bugs}, item 5). The specific
    concurrency levels chosen for the final sweep were derived partly from that calibration; the
    sweep design is defensible but not perfectly optimized as a result.
  \item \textbf{Three repetitions per cell} is a modest sample for characterizing the full
    variance at configurations near a capacity boundary, where we directly observed
    bimodal-looking outcomes (e.g., Shenandoah's spread even after noise control). A larger
    repetition count would better characterize this variance, particularly for any claim about
    typical versus worst-case behavior at the edge of capacity.
  \item \textbf{Spike-to-GC attribution is correlational} (\S\ref{sec:attribution}), not a proof
    that a given GC pause caused a given latency spike, particularly at high concurrency where
    general scheduling contention is also elevated.
\end{itemize}

\section{Conclusion}

The claim that garbage collection disqualifies Java from real-time agentic AI workloads is not
supported as a general statement: across three modern collectors and three heap sizes, GC
accounts for a small and stable share of tail latency whenever concurrency stays within a
reasonable margin of what the provisioned heap can sustain. But the claim is not simply false,
either --- and the more useful, less obvious finding of this study is that the breakdown which
does occur near a heap's true capacity limit is collector-specific in character, not just in
degree. G1, Generational ZGC, and Shenandoah each failed differently under the same pressure, and
the common advice to simply adopt a modern low-pause collector does not capture that difference.
The right choice depends on which failure mode a given system can least tolerate, and should be
verified against the system's actual workload shape rather than assumed from a collector's
general reputation --- the same discipline this report applied in finding that even its own test
environment's ambient load needed to be measured and controlled for, not assumed away.

\section{Reproducibility Appendix}

\begin{itemize}[leftmargin=1.5em]
  \item \textbf{Repository:} \code{gc-agentic-benchmark/} --- full harness source, build
    configuration (Maven, zero external dependencies), and orchestration scripts.
  \item \textbf{JDK requirement:} JDK 21+, and specifically a \textbf{native build matching your
    host CPU architecture} --- verify with \code{file \$(which java)} before trusting
    timing-sensitive results; see \S\ref{sec:jvm}.
  \item \textbf{Build:} \code{mvn package} from the project root.
  \item \textbf{Running a single calibration cell:} \code{scripts/pilot.sh} (fast, compressed
    timing --- see the caveat in \S\ref{sec:bugs}, item 5, before trusting its output as a
    precise predictor).
  \item \textbf{Running the full real study for one heap tier:}
    \code{JAVA\_BIN=<path to native JDK> ./scripts/run-real-study.sh <512m\textbar 1g\textbar 2g>}
    --- supports resuming a partially completed tier; any cell whose \code{status.txt} already
    exists is skipped.
  \item \textbf{Aggregating results:} \code{scripts/aggregate-real-results.sh} --- produces
    \code{data/results/aggregate.csv}, the source for every table and figure in this report.
  \item \textbf{Regenerating the figures:} \code{python3 scripts/generate\_charts.py} --- reads
    \code{aggregate.csv} directly and writes both SVGs in \code{figures/}; no external dependency
    (pure-stdlib SVG generation, since matplotlib was not available in the environment this was
    authored in).
  \item \textbf{Raw per-cell data:}
    \code{data/results/<heap>/<collector>/c<concurrency>/rep<N>/} contains the raw JFR recording,
    the raw per-token latency log, and a human-readable summary for every one of the 126
    (+27 re-verification) cells reported in this study.
  \item \textbf{Full narrative of the study's development} --- every bug found, every design
    decision and why it was made, is preserved in the project's development log alongside the
    repository.
\end{itemize}

\begin{thebibliography}{9}

\bibitem{infoq-trends-2025}
InfoQ Java Queue editorial team,
\emph{Java Trends Report 2025},
InfoQ, December 2025.
\url{https://www.infoq.com/articles/java-trends-report-2025/}

\bibitem{evans-2020}
B.~Evans,
\emph{Understanding Classic Java Garbage Collection},
InfoQ, May 2020.

\bibitem{marques-2025}
H.~Marques,
\emph{Java Concurrency from the Trenches: Lessons Learned in the Wild},
InfoQ Dev Summit Boston, 2025.

\end{thebibliography}

\end{document}
